\documentclass[10pt,a4paper]{article}
\usepackage[a4paper,margin=14mm,top=12mm,bottom=14mm]{geometry}
\usepackage{fontspec}
\setmainfont{Liberation Sans}
\usepackage{polyglossia}
\setdefaultlanguage{french}
\usepackage{amsmath}
\usepackage{xcolor}
\usepackage{graphicx}
\usepackage[most]{tcolorbox}
\usepackage{tabularx,booktabs,array,multirow}
\usepackage{enumitem}
\usepackage{fancyhdr}
\usepackage{lastpage}
\usepackage[hidelinks]{hyperref}
\definecolor{uimm}{HTML}{005677}
\definecolor{cyan}{HTML}{13A2DD}
\definecolor{fond}{HTML}{E6F2FA}
\definecolor{ok}{HTML}{2E7D32}
\definecolor{ko}{HTML}{B3261E}
\definecolor{violet}{HTML}{7B1FA2}
\setlength{\parindent}{0pt}
\setlength{\parskip}{3pt}
\setlist{nosep}
\setlist[itemize]{leftmargin=13pt,label=\textcolor{cyan}{\textbullet}}
\newcommand{\titre}[1]{\vspace{5pt}{\color{uimm}\large\bfseries #1}\par\vspace{-3pt}{\color{cyan}\rule{\linewidth}{1.2pt}}\par\vspace{0pt}}
\tcbset{boite/.style={colback=fond,colframe=uimm,boxrule=0.6pt,arc=2pt,left=6pt,right=6pt,top=4pt,bottom=4pt},
        defi/.style={colback=violet!6,colframe=violet,boxrule=0.8pt,arc=2pt,left=6pt,right=6pt,top=4pt,bottom=4pt,
                     title={\bfseries Défi},coltitle=white,colbacktitle=violet,fonttitle=\small}}
\newcommand{\bandeau}[3]{%
\begin{tcolorbox}[colback=uimm,colframe=uimm,arc=2pt,left=4pt,right=8pt,top=4pt,bottom=4pt]
\begin{minipage}[c]{18mm}\setlength{\fboxsep}{1.2pt}\fcolorbox{white}{white}{\includegraphics[height=16.5mm]{logo_uimm.jpg}}\end{minipage}\hspace{4mm}%
\begin{minipage}[c]{\dimexpr\linewidth-22mm\relax}\color{white}
{\small #1}\\[1pt]{\LARGE\bfseries #2}\\[2pt]{\large #3}\hfill{\small CQPM Technicien qualité}
\end{minipage}
\end{tcolorbox}}
\newcommand{\reponse}[1][1]{\par\vspace{2pt}{\color{gray!60}\foreach \i in {1,...,#1}{\rule{\linewidth}{0.3pt}\par\vspace{9pt}}}}
\usepackage{pgffor}
\pagestyle{fancy}\fancyhf{}
\renewcommand{\headrulewidth}{0pt}
\newcommand{\pied}[1]{\fancyfoot[L]{\footnotesize\color{uimm}#1}\fancyfoot[R]{\footnotesize\color{uimm}Page \thepage{} sur \pageref{LastPage}}}

\usepackage{tikz}
\usepackage{longtable}
\renewcommand{\bandeau}[3]{%
\begin{tcolorbox}[colback=uimm,colframe=uimm,arc=2pt,left=4pt,right=8pt,top=4pt,bottom=4pt]
\begin{minipage}[c]{18mm}\setlength{\fboxsep}{1.2pt}\fcolorbox{white}{white}{\includegraphics[height=16.5mm]{logo_uimm.jpg}}\end{minipage}\hspace{4mm}%
\begin{minipage}[c]{\dimexpr\linewidth-22mm\relax}\color{white}
{\small #1}\\[1pt]{\LARGE\bfseries #2}\\[2pt]{\large #3}
\end{minipage}
\end{tcolorbox}}
% en-tête d'une fiche : code, titre, question, bloc, durée, organisation, niveau
\newcommand{\entete}[7]{%
\begin{tcolorbox}[colback=uimm,colframe=uimm,arc=2pt,left=4pt,right=6pt,top=4pt,bottom=4pt]
\begin{minipage}[c]{18mm}\setlength{\fboxsep}{1.2pt}\fcolorbox{white}{white}{\includegraphics[height=16.5mm]{logo_uimm.jpg}}\end{minipage}\hspace{3mm}%
\begin{minipage}[c]{\dimexpr\linewidth-58mm\relax}\color{white}\raggedright
{\small Base SPC · #4}\\[1pt]{\Large\bfseries #1 · #2}\\[2pt]{\itshape #3}
\end{minipage}\hfill
\begin{minipage}[c]{33mm}\color{white}\small\raggedleft
Durée : \textbf{#5}\\ \textbf{#6}\\ Niveau #7 sur 3
\end{minipage}
\end{tcolorbox}}
\newcommand{\tagE}{{\setlength{\fboxsep}{1.5pt}\colorbox{cyan}{\color{white}\scriptsize\bfseries Ellistat}}}
\newcommand{\tagW}{{\setlength{\fboxsep}{1.5pt}\colorbox{uimm}{\color{white}\scriptsize\bfseries Page web}}}
\tcbset{pred/.style={colback=orange!6,colframe=orange!80!black,boxrule=0.6pt,arc=2pt,left=6pt,right=6pt,top=3pt,bottom=1pt,
                    title={\bfseries Je prédis, avant de commencer},coltitle=white,colbacktitle=orange!80!black,fonttitle=\small},
        retenir/.style={colback=fond,colframe=uimm,boxrule=0.6pt,arc=2pt,left=6pt,right=6pt,top=3pt,bottom=3pt,
                    title={\bfseries À retenir},coltitle=white,colbacktitle=uimm,fonttitle=\small},
        plus/.style={colback=violet!6,colframe=violet,boxrule=0.6pt,arc=2pt,left=6pt,right=6pt,top=3pt,bottom=3pt,
                    title={\bfseries Pour aller plus loin},coltitle=white,colbacktitle=violet,fonttitle=\small},
        act/.style={colback=uimm!10,colframe=uimm,boxrule=0.5pt,arc=2pt,left=5pt,right=5pt,top=2pt,bottom=2pt,before skip=7pt,after skip=3pt},
        verif/.style={colback=orange!6,colframe=orange!80!black,boxrule=0.6pt,arc=2pt,left=6pt,right=6pt,top=3pt,bottom=3pt,
                    title={\bfseries À vérifier sur poste avant la première séance},coltitle=white,colbacktitle=orange!80!black,fonttitle=\small}}
\newcommand{\fichier}[1]{\texttt{\small #1}}
\newcommand{\menu}[1]{\textsf{\textbf{#1}}}
\newcommand{\champ}[1]{\rule{#1}{0.3pt}}

\pied{Base SPC · guide formateur}
\begin{document}
\thispagestyle{fancy}
\bandeau{Guide formateur · ne pas distribuer}{Base d'activités SPC}{Mode d'emploi, organisation, corrigés des 36 fiches}

\titre{1. La base en une page}
\begin{itemize}
\item \textbf{36 activités, environ 19 heures}, en quatre blocs : E (prise en main d'Ellistat), M (moyen de mesure), A (cartes aux attributs), C (cartes aux mesures et capabilité). On choisit, on ne fait pas tout.
\item \textbf{Une fiche par activité}, toujours bâtie de la même façon : le but, ce qu'il faut, « je prédis », « à faire », « j'observe, je décide », « à retenir ». Une seule question par fiche.
\item \textbf{Prédire, faire, comparer.} Le stagiaire écrit ce qu'il attend avant de mesurer ou de calculer. C'est ce qui ralentit et qui ancre.
\item \textbf{Ils mesurent, Ellistat calcule.} Aucune fiche ne demande Excel. Les données de secours se collent depuis la page « Données à copier ».
\item \textbf{Le catalogue} (\fichier{index.html}) : on y règle le nombre de stagiaires et le temps disponible, on coche les activités, il affiche l'organisation et le matériel.
\item \textbf{Tout est dans le dossier.} Les pages web, les fiches et les données ne renvoient à aucun autre dossier. Seuls le guide en ligne d'Ellistat et le livret Bosch sont cités comme sources externes.
\end{itemize}

\titre{2. Selon le nombre de stagiaires}
\begin{itemize}
\item \textbf{Binôme par défaut} pour les mesures : l'un mesure, l'autre note, puis on échange. Effectif impair : un trinôme (le troisième chronomètre ou saisit).
\item \textbf{Un seul stagiaire} : toutes les fiches en binôme se font seul, le formateur note. M9 demande deux opérateurs : le formateur est le second.
\item \textbf{Matériel par poste} : un micromètre et une cale par binôme pour M2 à M7 et M9 ; cinq cales par binôme pour M8 ; un coffret pour M10.
\item \textbf{Activités de groupe} : M11 (R\&R) demande 3 opérateurs ; au-delà, les autres notent. M12 (Trois cales) fonctionne de 3 à 8 stagiaires. M13 se fait à partir de 2.
\item \textbf{Ellistat} : un compte par stagiaire pour le bloc E. Ensuite, un compte par binôme suffit.
\end{itemize}

\titre{3. Parcours suggérés}
\renewcommand{\arraystretch}{1.25}
\begin{tabularx}{\linewidth}{@{}p{43mm} X r@{}}
\toprule
\textbf{Parcours} & \textbf{Activités, dans l'ordre} & \textbf{Durée}\\\midrule
Journée très pratique, public d'opérateurs (7 h) & E1, E2, E3, E4 · M1, M2, M3, M4, M5, M7, M9 · A1, A2, A3, A4, A6 & 6 h 15\\
Découverte d'Ellistat & E1 à E6 & 1 h 30\\
Moyen de mesure, une journée & E1, E2, E3 · M1 à M10 · M11 si le groupe avance & 5 h 20 à 6 h 30\\
Cartes aux attributs, une demi-journée & A1 à A6, A8, A9 & 3 h 20\\
Cartes aux mesures et capabilité, une journée & C1, C2, C3 · E5 · C4, C6, C7 & 4 h 20\\
Groupe solide, synthèse & C5, M8, M12 ou M13, C8 & 3 h 20 à 4 h 15\\
\bottomrule
\end{tabularx}
{\small Les durées n'incluent ni les pauses ni les bilans. Avec un public d'exécutants : niveaux 1 et 2 seulement, et les activités de réserve sont M6 et M10 (très concrètes).}

\titre{4. Ellistat : ce qu'on sait, ce qu'il faut vérifier}
\begin{itemize}
\item \textbf{Sous-groupes} : une ligne de la grille = un sous-groupe ; le champ « Répétition » de la première colonne donne le nombre de colonnes du sous-groupe (réponse d'Ellistat du 10/11/2026).
\item \textbf{Cartes et capabilité} (établi sur une copie d'écran du 10/11/2026, recalculé en Python) : écart-type court terme = écart-type groupé ; limites à $\bar{\bar X} \pm 3\sigma/\sqrt n$ ; Pp et Ppk avec l'écart-type global ; Rs = Pp / Cp ; Rr = Ppk / Pp.
\item \textbf{Justesse} (guide Ellistat, page mise à jour le 06/11/2026) : au moins 25 mesures visées, 10 au minimum ; voyants : résolution verte jusqu'à T/20, orange jusqu'à T/10 ; Cg et Cgk verts au-dessus de 1,33.
\end{itemize}
\begin{tcolorbox}[verif]
\begin{itemize}
\item Le séparateur décimal accepté par la grille (activité E2).
\item L'écran \menu{Justesse} et l'écran des cartes aux attributs : le guide n'en donne pas le pas à pas.
\item Conventions de calcul non confirmées par Ellistat : Cg en 6\,s ou en 4\,s ; R\&R avec le facteur 6 ou 5,15 ; regroupement de l'interaction. Fiche de test : \fichier{Ellistat\_tests\_conventions.pdf}.
\item Les corrigés chiffrés de Cg, Cgk et \%GRR de ce guide suivent les formules du livret Bosch. Ils sont à comparer à l'écran.
\item Le comportement avec Répétition = 1, le nombre de points des règles d'alerte, la carte de Laney.
\end{itemize}
\end{tcolorbox}

\titre{5. La page « Tirer un lot » : les scénarios}
\renewcommand{\arraystretch}{1.2}
\begin{tabularx}{\linewidth}{@{}c X X@{}}
\toprule
\textbf{Scénario} & \textbf{Pièces défectueuses} & \textbf{Défauts par unité}\\\midrule
A & stable à 4 \% & stable à 3 défauts\\
B & 4 \% puis 10 \% à partir du lot 21 & 3 puis 6 défauts à partir du n° 21\\
C & 4 \%, puis dérive de 0,4 point par lot à partir du lot 16 & 3, puis dérive de 0,2 défaut par unité à partir du n° 16\\
D & 4 \% puis 1,5 \% à partir du lot 21 (amélioration) & 3 puis 1,2 défaut à partir du n° 21\\
E & 4 \%, sauf le lot 18 à 15 \% & 3, sauf le n° 18 à 9 défauts\\
\bottomrule
\end{tabularx}
{\small Le même scénario et le même numéro de poste donnent toujours les mêmes tirages : donner un numéro différent à chaque binôme. Le bouton « Révéler le taux caché » est en bas de la page.}

\titre{6. Mesurer avec des cales : les précautions}
\begin{itemize}
\item Cales et instruments dans la salle au moins une heure avant. Les cales sont définies à 20 °C.
\item Tenir les cales par les bords, avec des gants ou un chiffon. Ne pas les garder dans la main (activité M6).
\item Faces propres. Une face oxydée ou rayée n'adhère plus : la cale ne sert plus à composer une valeur (M10).
\item L'étiquette d'un coffret (par exemple « DIN 861/I ») indique la classe d'origine. Elle ne remplace pas un certificat d'étalonnage. Sans certificat, on prend la valeur nominale et on le dit.
\item Micromètre : zéro contrôlé, cliquet, touches propres. Un micromètre 25-50 se règle sur sa cale de 25 mm.
\end{itemize}

\titre{7. Toutes les activités}
{\small\renewcommand{\arraystretch}{1.12}\begin{longtable}{@{}l p{74mm} r l c c c@{}}\toprule
\textbf{Code} & \textbf{Activité} & \textbf{Durée} & \textbf{Organisation} & \textbf{Niveau} & \textbf{Ellistat} & \textbf{Mesures}\\\midrule\endhead
\multicolumn{7}{@{}l}{\textbf{\color{uimm}Bloc E · Prise en main d'Ellistat} \quad (90 min)}\\
E1 & Mon premier projet & 10 min & seul & 1 & oui & \\
E2 & Saisir, coller, corriger & 15 min & seul & 1 & oui & \\
E3 & Régler une colonne & 15 min & seul & 1 & oui & \\
E4 & Mon premier histogramme & 20 min & seul & 1 & oui & \\
E5 & Une ligne, un sous-groupe & 20 min & seul & 2 & oui & \\
E6 & Garder une trace & 10 min & seul & 1 & oui & \\
\multicolumn{7}{@{}l}{\textbf{\color{uimm}Bloc M · Le moyen de mesure} \quad (495 min)}\\
M1 & La résolution d'abord & 15 min & en binôme & 1 &  & oui\\
M2 & Dix mesures, une cale & 20 min & en binôme & 1 &  & oui\\
M3 & Mon micromètre est-il capable ? & 40 min & en binôme & 2 & oui & oui\\
M4 & Pied à coulisse contre micromètre & 30 min & en binôme & 2 & oui & oui\\
M5 & Avec ou sans cliquet & 30 min & en binôme & 2 & oui & oui\\
M6 & La cale dans la main & 25 min & en binôme & 2 &  & oui\\
M7 & Jusqu'où peut-il aller ? & 20 min & en binôme & 2 & oui & \\
M8 & Juste partout ? & 45 min & en binôme & 3 & oui & oui\\
M9 & Deux opérateurs, une cale & 30 min & en binôme & 2 &  & oui\\
M10 & Fabriquer son étalon & 25 min & en binôme & 2 & oui & oui\\
M11 & L'opérateur compte-t-il ? & 70 min & en groupe & 3 & oui & oui\\
M12 & Trois cales, trois questions & 100 min & rotation sur postes & 3 &  & oui\\
M13 & Bon ou mauvais : sommes-nous d'accord ? & 45 min & en groupe & 3 & oui & oui\\
\multicolumn{7}{@{}l}{\textbf{\color{uimm}Bloc A · Cartes aux attributs} \quad (225 min)}\\
A1 & Défectueuse ou défauts ? & 15 min & en binôme & 1 &  & \\
A2 & Ma première carte p, à la main & 35 min & en binôme & 1 &  & \\
A3 & La même carte dans Ellistat & 20 min & en binôme & 1 & oui & \\
A4 & Le procédé se dérègle & 30 min & en binôme & 2 & oui & \\
A5 & p ou np ? Et si les lots changent de taille ? & 25 min & en binôme & 2 & oui & \\
A6 & Compter les défauts : la carte c & 30 min & en binôme & 2 & oui & \\
A7 & Des unités de taille variable : la carte u & 25 min & en binôme & 3 & oui & \\
A8 & Quelle période de référence ? & 25 min & en binôme & 2 & oui & \\
A9 & La taille du lot compte & 20 min & en binôme & 2 &  & \\
\multicolumn{7}{@{}l}{\textbf{\color{uimm}Bloc C · Cartes aux mesures et capabilité} \quad (350 min)}\\
C1 & La carte sonne-t-elle ? & 30 min & en binôme & 1 &  & \\
C2 & La machine est-elle capable ? & 30 min & en binôme & 1 &  & \\
C3 & La voiture et le garage & 30 min & en binôme & 2 &  & \\
C4 & Ma carte, à la main puis dans Ellistat & 60 min & en binôme & 2 & oui & \\
C5 & Le mauvais découpage & 20 min & en binôme & 3 & oui & \\
C6 & Cinquante butées qui dérivent & 30 min & en binôme & 2 & oui & \\
C7 & Trois enquêtes & 60 min & en binôme & 2 & oui & \\
C8 & Du moyen de mesure à la décision & 90 min & en binôme & 3 & oui & \\
\bottomrule\end{longtable}}

\newpage
\titre{8. Préparation, animation et corrigés, fiche par fiche}
{\small Chaque corrigé porte le code de sa fiche ; « Q1 », « Q2 »… renvoient aux questions de la rubrique « J'observe, je décide » de cette fiche. « Résultats propres à chaque poste » : les valeurs dépendent des mesures du jour ; le guide donne alors les valeurs du jeu de secours.}
\section*{Bloc E · Prise en main d'Ellistat}
\begin{tcolorbox}[act]\textbf{E1 · Mon premier projet} \hfill {\small 10 min · seul · niveau 1 · Ellistat}\end{tcolorbox}
\textbf{Préparation.} Vérifier la veille que chaque compte ouvre bien la tuile Data Analysis (un module absent de l'accueil vient de la licence).\par
\textbf{Animation.} Faire le premier projet ensemble, écran projeté, puis laisser chacun refaire seul. Convention de nommage : « Thème Nom », pour retrouver les projets de chacun.\par
\textbf{\color{ok}Corrigé de la fiche E1.} \textbf{Prédiction} le projet. \textbf{Q1} projet ; grilles ; grille. \textbf{Q2} Gage R\&R, Graphiques, Statistique descriptive, Statistiques inférentielles… Pour l'instrument : Gage R\&R. \textbf{Q3} Aucun projet n'est ouvert : en choisir un dans le sélecteur.\par
{\small\textbf{Repli ou remarque.} Le guide Ellistat ne décrit pas l'écran d'accueil en détail : vérifier les libellés sur votre compte.}\par
\vspace{3pt}
\begin{tcolorbox}[act]\textbf{E2 · Saisir, coller, corriger} \hfill {\small 15 min · seul · niveau 1 · Ellistat}\end{tcolorbox}
\textbf{Préparation.} Tester soi-même la virgule et le point : le guide Ellistat ne dit pas lequel est accepté.\par
\textbf{Animation.} Laisser les stagiaires découvrir le séparateur : c'est une vraie question, pas un piège. Faire noter la réponse au tableau : elle sert pour toute la journée.\par
\textbf{\color{ok}Corrigé de la fiche E2.} \textbf{Q1} À constater sur poste et à noter. \textbf{Q2} Par copier-coller dans la grille : l'import crée une grille, il ne complète pas. \textbf{Q3} Une ligne vide ou une valeur mal lue fausse tous les calculs sans prévenir.\par
{\small\textbf{Repli ou remarque.} Si le collage échoue, passer par \menu{Importer depuis Excel}.}\par
\vspace{3pt}
\begin{tcolorbox}[act]\textbf{E3 · Régler une colonne} \hfill {\small 15 min · seul · niveau 1 · Ellistat}\end{tcolorbox}
\textbf{Animation.} Insister sur le panneau de colonne : tout se règle là, y compris plus tard les sous-groupes.\par
\textbf{\color{ok}Corrigé de la fiche E3.} \textbf{Prédiction} non : la tolérance vient du plan, pas des mesures. \textbf{Q1} Moyenne 11,887 ; écart-type 0,026 (30 valeurs). Elle est sous le milieu de la tolérance. \textbf{Q2} Une capabilité compare la dispersion à la tolérance : sans tolérance, rien à comparer. \textbf{Q3} À revenir en arrière sans rien perdre.\par
{\small\textbf{Repli ou remarque.} Si l'onglet \menu{Stats} n'affiche pas l'écart-type, le lire dans l'onglet \menu{données}.}\par
\vspace{3pt}
\begin{tcolorbox}[act]\textbf{E4 · Mon premier histogramme} \hfill {\small 20 min · seul · niveau 1 · Ellistat}\end{tcolorbox}
\textbf{Préparation.} Repérer sur votre compte la boîte de dialogue de \menu{Ajouter des valeurs aléatoires} : le guide ne la décrit pas.\par
\textbf{Animation.} Faire répondre à la prédiction avant d'ouvrir l'écran. Relier à l'image de la voiture et du garage (simulateur de capabilité, activité C3).\par
\textbf{\color{ok}Corrigé de la fiche E4.} \textbf{Q1} 2 pièces sous 11,85 (11,84 deux fois), aucune au-dessus de 12,00. \textbf{Q2} La population est estimée avec la cloche : elle prolonge ce qu'on a vu à toutes les pièces futures. Avec la moyenne 11,887 et l'écart-type 0,026, la loi normale donne environ 8 \% sous 11,85. \textbf{Q3} Non, elle est décalée vers la limite basse : le risque est en bas.\par
{\small\textbf{Repli ou remarque.} Vérifier à l'écran le pourcentage annoncé par Ellistat et le noter dans ce guide.}\par
\vspace{3pt}
\begin{tcolorbox}[act]\textbf{E5 · Une ligne, un sous-groupe} \hfill {\small 20 min · seul · niveau 2 · Ellistat}\end{tcolorbox}
\textbf{Animation.} C'est la clé de toutes les cartes aux mesures : prendre le temps. Le mot « Répétition » est trompeur, Ellistat le reconnaît (réponse du 10/11/2026).\par
\textbf{\color{ok}Corrigé de la fiche E5.} \textbf{Prédiction} 6 points. \textbf{Q1} Une ligne = un prélèvement de 5 pièces ; une colonne = le rang de la pièce dans le prélèvement. \textbf{Q2} 6 points, moyenne générale 11,8867. Avec Répétition = 1, P1 redevient une caractéristique seule (6 valeurs individuelles ; comportement exact à confirmer à l'écran). \textbf{Q3} Le calcul est le même : moyenne et dispersion à l'intérieur de la ligne. Ce qui compte : les pièces ont été prélevées ensemble.\par
{\small\textbf{Repli ou remarque.} Un fichier avec des lignes vides crée des sous-groupes vides : toujours vérifier le nombre de lignes.}\par
\vspace{3pt}
\begin{tcolorbox}[act]\textbf{E6 · Garder une trace} \hfill {\small 10 min · seul · niveau 1 · Ellistat}\end{tcolorbox}
\textbf{Animation.} À placer en fin de séance : chacun repart avec un graphique nommé.\par
\textbf{\color{ok}Corrigé de la fiche E6.} \textbf{Q1} Non : c'est une image figée. \textbf{Q2} Un titre qui dit quoi, sur quoi et quand : « Histogramme, 30 butées, 25/11 ».\par
{\small\textbf{Repli ou remarque.} Un projet « public » serait visible de tous les comptes de l'organisme : préférer les utilisateurs autorisés.}\par
\vspace{3pt}
\section*{Bloc M · Le moyen de mesure}
\begin{tcolorbox}[act]\textbf{M1 · La résolution d'abord} \hfill {\small 15 min · en binôme · niveau 1 · mesures}\end{tcolorbox}
\textbf{Préparation.} Relever la résolution réelle de vos instruments : micromètre au 0,01 ou au 0,001 selon le modèle.\par
\textbf{Animation.} Aucun calcul de capabilité ici : une division. C'est la porte d'entrée de tout le bloc. Faire toucher les instruments avant de calculer.\par
\textbf{\color{ok}Corrigé de la fiche M1.} Réglet (0,5 mm) : 333 \% ; 1000 \% ; 5000 \%. Pied à coulisse (0,02) : 13,3 \% ; 40 \% ; 200 \%. Micromètre (0,001) : 0,67 \% ; 2 \% ; 10 \%. Avec un micromètre au 0,01 : 6,7 \% ; 20 \% ; 100 \%.\par \textbf{Q1} Le micromètre au 0,001 pour la butée ; aucun des trois pour 0,01 mm (il faudrait un comparateur ou une colonne de mesure). \textbf{Q2} Non : il est aveugle à la dispersion.\par
{\small\textbf{Repli ou remarque.} Critère : livret Bosch n° 10, fig. 3 (résolution au plus 5 \% de T). Ellistat applique le même seuil (voyant vert jusqu'à T/20).}\par
\vspace{3pt}
\begin{tcolorbox}[act]\textbf{M2 · Dix mesures, une cale} \hfill {\small 20 min · en binôme · niveau 1 · mesures}\end{tcolorbox}
\textbf{Préparation.} Cales et micromètres dans la salle une heure avant. Zéro des micromètres contrôlé.\par
\textbf{Animation.} Ne pas parler de Cg : seulement « groupé » et « centré ». Dessiner au tableau les quatre cibles (groupé ou non, centré ou non).\par
\textbf{\color{ok}Corrigé de la fiche M2.} Résultats propres à chaque poste. Ordres de grandeur avec un micromètre au 0,001 en bon état : étendue de 1 à 3 µm, écart à la cale de 0 à 2 µm.\par \textbf{Q3} Poser la pièce fait partie de la mesure : sans cela on sous-estime la dispersion (livret Bosch, § 4.1).\par
{\small\textbf{Repli ou remarque.} Si toutes les valeurs sont identiques : relier à M1 (résolution) ou vérifier que la cale a bien été retirée.}\par
\vspace{3pt}
\begin{tcolorbox}[act]\textbf{M3 · Mon micromètre est-il capable ?} \hfill {\small 40 min · en binôme · niveau 2 · Ellistat · mesures}\end{tcolorbox}
\textbf{Préparation.} Dérouler une étude de justesse sur votre compte : le guide Ellistat ne donne pas le pas à pas de cet écran. La cale de 10 mm n'est pas au milieu de la tolérance (11,925) : le dire. Pour un étalon à 12 mm, voir M10.\par
\textbf{Animation.} 25 mesures prennent 10 à 15 minutes : le secrétaire saisit au fur et à mesure. Avec $T = 0{,}15$, le micromètre sera capable. L'intérêt vient avec M4, M5 et M7.\par
\textbf{\color{ok}Corrigé de la fiche M3.} Résultats propres à chaque poste. Sur le jeu de secours (colonne Micrometre, cale 10,000) : moyenne 10,0004 ; écart-type 0,00081 ; avec les formules du livret Bosch, Cg = 6,17 et Cgk = 6,02.\par \textbf{Q3} Cgk est toujours inférieur ou égal à Cg : il retire le décalage. Les deux sont égaux seulement si la moyenne tombe sur la valeur de la cale.\par \textbf{Pour aller plus loin} La page applique Cg = 0,2 T / (6 s). Si Ellistat affiche 1,5 fois plus, il compte la dispersion en 4 s : le noter (point non confirmé par Ellistat).\par
{\small\textbf{Repli ou remarque.} Sans mesures : jeu « Secours : micromètre et pied à coulisse », colonne Micrometre.}\par
\vspace{3pt}
\begin{tcolorbox}[act]\textbf{M4 · Pied à coulisse contre micromètre} \hfill {\small 30 min · en binôme · niveau 2 · Ellistat · mesures}\end{tcolorbox}
\textbf{Animation.} Faire annoncer la prédiction avant de mesurer. Le pied à coulisse donne souvent 2 ou 3 valeurs seulement (10,00 ; 10,02) : c'est l'effet de la résolution, à faire constater sur le graphique.\par
\textbf{\color{ok}Corrigé de la fiche M4.} Jeu de secours (25 mesures chacun) : micromètre $s = 0{,}0008$, Cg = 6,17, Cgk = 6,02 ; pied à coulisse $s = 0{,}0099$, Cg = 0,51, Cgk = 0,45.\par \textbf{Q2} La résolution : 0,02 / 0,15 = 13,3 \%, au-dessus de 10 \% (voyant rouge). \textbf{Q3} Environ 12 fois sur le jeu de secours (F = 148, seuil 2,27 : différence nette). L'écart entre les moyennes (+1,2 µm) n'est pas démontré (t = 0,63, seuil 2,06).\par \textbf{Q4} Non. Résolution 0,02 au plus 5 \% de T : il convient à partir de $T = 0{,}4$ mm.\par
{\small\textbf{Repli ou remarque.} Jeu « Secours : micromètre et pied à coulisse ».}\par
\vspace{3pt}
\begin{tcolorbox}[act]\textbf{M5 · Avec ou sans cliquet} \hfill {\small 30 min · en binôme · niveau 2 · Ellistat · mesures}\end{tcolorbox}
\textbf{Préparation.} Sans cliquet, ne pas forcer : un serrage excessif abîme le micromètre et la cale.\par
\textbf{Animation.} Activité très parlante pour des opérateurs : elle part de leur geste. Faire écrire la consigne (Q3) et en afficher une au tableau.\par
\textbf{\color{ok}Corrigé de la fiche M5.} Résultats propres à chaque poste. Jeu de secours : avec cliquet $s = 0{,}0005$, biais +0,0003, Cg = 10,9 ; sans cliquet $s = 0{,}0023$, biais $-0{,}0032$, Cg = 2,1 et Cgk = 1,7.\par Écart des moyennes $-3{,}5$ µm, réel (t = 5,6, seuil 2,1). Dispersion 5 fois plus grande (F = 26, seuil 3,0). Un serrage à la main écrase davantage : lectures plus petites et moins répétables. Le sens peut différer selon l'opérateur.\par \textbf{Q3} Par exemple : « Mesurer au cliquet, trois crans, pièce posée bien à plat sur les touches. »\par
{\small\textbf{Repli ou remarque.} Jeu « Secours : avec et sans cliquet ».}\par
\vspace{3pt}
\begin{tcolorbox}[act]\textbf{M6 · La cale dans la main} \hfill {\small 25 min · en binôme · niveau 2 · mesures}\end{tcolorbox}
\textbf{Préparation.} Prendre une cale longue : l'effet est proportionnel à la longueur. Une cale de 5 mm ne montre rien.\par
\textbf{Animation.} Le résultat surprend toujours : faire parier avant. Relier aux mesures de la journée : c'est pour cela qu'on tient les cales par les bords.\par
\textbf{\color{ok}Corrigé de la fiche M6.} \textbf{Q1} Propre à chaque poste. Jeu de secours : +3 µm au départ, retour en 3 minutes environ. \textbf{Q2} $0{,}025 \times 11{,}5 \times 10 = 2{,}9$ µm. \textbf{Q3} Stabiliser pièces et instruments dans la salle, tenir par les bords ou avec des gants, ne pas mesurer une pièce qui sort d'usinage.\par Le coefficient de 11,5 µm par mètre et par degré est une valeur usuelle pour l'acier : à confirmer sur la documentation du fabricant des cales. Les cales sont définies à 20 °C.\par
{\small\textbf{Repli ou remarque.} Jeu « Secours : cale de 25 mm réchauffée ».}\par
\vspace{3pt}
\begin{tcolorbox}[act]\textbf{M7 · Jusqu'où peut-il aller ?} \hfill {\small 20 min · en binôme · niveau 2 · Ellistat}\end{tcolorbox}
\textbf{Animation.} Aucune mesure nouvelle : l'activité se fait sur les données de M3. Bonne activité d'attente pour un binôme en avance.\par
\textbf{\color{ok}Corrigé de la fiche M7.} Jeu de secours (micromètre, $s = 0{,}00081$, biais +0,00036), formules du livret Bosch : $T = 0{,}15$ : Cg 6,17, Cgk 6,02 ; $T = 0{,}05$ : 2,06 et 1,91 ; $T = 0{,}02$ : 0,82 et 0,67 ; $T = 0{,}01$ : 0,41 et 0,26.\par \textbf{Q1} 0,05 mm. \textbf{Q2} $T_{\min} = 0{,}0036 + 0{,}0323 = 0{,}036$ mm. \textbf{Q3} Cg est divisé par 3.\par La formule de $T_{\min}$ vient des deux conditions Cg et Cgk au moins égaux à 1,33. Si Ellistat annonce une autre valeur, sa convention diffère (point non confirmé).\par
{\small\textbf{Repli ou remarque.} Jeu « Secours : micromètre et pied à coulisse », colonne Micrometre.}\par
\vspace{3pt}
\begin{tcolorbox}[act]\textbf{M8 · Juste partout ?} \hfill {\small 45 min · en binôme · niveau 3 · Ellistat · mesures}\end{tcolorbox}
\textbf{Préparation.} Cinq cales par poste : avec plusieurs coffrets, chaque binôme a les siennes. Sinon, faire tourner les cales entre les postes.\par
\textbf{Animation.} 50 mesures : prévoir 25 minutes de mesure. Répartir : un mesure, un note, on échange à mi-parcours. Faire tracer le biais en fonction de la longueur avant d'ouvrir Ellistat.\par
\textbf{\color{ok}Corrigé de la fiche M8.} Résultats propres à chaque poste. Jeu de secours (simulé) : biais moyens $-2{,}7$ ; $-1{,}0$ ; $+1{,}3$ ; $+3{,}2$ ; $+3{,}7$ µm. Pente 0,00034 mm par mm, très significative ; le biais s'annule vers 12,4 mm.\par \textbf{Q3} Constant : régler le zéro ou étalonner. Variable : réduire la plage d'utilisation, étalonner près de la cote mesurée, ou changer d'instrument.\par Le livret Bosch (§ 4.4) propose de répéter l'étude de justesse sur plusieurs étalons, au moins 5 si possible.\par
{\small\textbf{Repli ou remarque.} Jeu « Linéarité : 5 cales × 10 mesures (simulé) », bouton « Copier les 10 mesures » de chaque ligne.}\par
\vspace{3pt}
\begin{tcolorbox}[act]\textbf{M9 · Deux opérateurs, une cale} \hfill {\small 30 min · en binôme · niveau 2 · mesures}\end{tcolorbox}
\textbf{Préparation.} Avec un nombre impair de stagiaires, former un trinôme : trois opérateurs, comparés deux à deux.\par
\textbf{Animation.} Personne n'a « raison » : la cale dit qui est le plus juste, pas qui est le meilleur opérateur. Prépare directement le R\&R (M12).\par
\textbf{\color{ok}Corrigé de la fiche M9.} Résultats propres à chaque poste. Jeu de secours : A moyenne 10,0001, B moyenne 10,0011 ; écart +1,1 µm, réel (t = 4,7, seuil 2,05). Dispersions comparables (F = 1,16, seuil 2,98).\par \textbf{Q3} Force de serrage, nombre de crans, position de la cale sur les touches, lecture. La valeur de la cale tranche : celui dont la moyenne en est la plus proche est le plus juste.\par
{\small\textbf{Repli ou remarque.} Jeu « Secours : deux opérateurs ».}\par
\vspace{3pt}
\begin{tcolorbox}[act]\textbf{M10 · Fabriquer son étalon} \hfill {\small 25 min · en binôme · niveau 2 · Ellistat · mesures}\end{tcolorbox}
\textbf{Préparation.} Vérifier dans le coffret la présence des cales 10 ; 2 ; 1,35 et 0,5 mm, et l'état des faces : une face oxydée ou rayée n'adhère plus. Montrer le geste d'adhérence avant de laisser faire.\par
\textbf{Animation.} Geste de métrologue, gratifiant. Ne jamais séparer deux cales en tirant. Le livret Bosch demande un étalon si possible au milieu de la tolérance (§ 4.1) et, avec deux étalons, aux deux limites (§ 4.4).\par
\textbf{\color{ok}Corrigé de la fiche M10.} Résultats propres à chaque poste. La pile vaut la somme des valeurs des deux cales. En classe 1, chaque cale est garantie à quelques dixièmes de micromètre : l'écart possible reste sous la résolution du micromètre.\par \textbf{Q2} L'instrument peut être juste à 10 mm et décalé à 12 mm (voir M8). On le qualifie à la longueur où il sert.\par Milieu de la tolérance : 11,925 = 9,5 + 1,42 + 1,005, si le coffret contient ces cales.\par
{\small\textbf{Repli ou remarque.} Si les cales n'adhèrent pas : les nettoyer de nouveau, sinon utiliser la cale de 10 mm seule et le dire.}\par
\vspace{3pt}
\begin{tcolorbox}[act]\textbf{M11 · L'opérateur compte-t-il ?} \hfill {\small 70 min · en groupe · niveau 3 · Ellistat · mesures}\end{tcolorbox}
\textbf{Préparation.} Choisir la veille 10 butées étalées de 11,85 à 12,00, les numéroter, tracer le trait repère. Il faut au moins 3 personnes pour mesurer. À 4 ou 5, les autres notent. À 2, le formateur est le troisième opérateur.\par
\textbf{Animation.} Corrigé détaillé : guide « L'opérateur compte-t-il ? » (fichier \fichier{Guide\_RR\_formateur.pdf}). Calibrage de la page : bouton « Exemple du livret Bosch », \%GRR = 17,95 \%, ndc = 15.\par
\textbf{\color{ok}Corrigé de la fiche M11.} \textbf{Q2} La page applique la méthode du livret Bosch (facteur 6, interaction regroupée si elle n'est pas significative). Ellistat peut utiliser d'autres conventions : point non confirmé, à tester avec l'exemple du livret (\fichier{RR\_Bosch\_p14.xlsx}).\par Seuils du livret Bosch (§ 4.2) : \%GRR au plus 10 \% capable ; de 10 à 30 \% capable sous conditions ; plus de 30 \% non capable.\par
{\small\textbf{Repli ou remarque.} Sans butées : exemple du livret Bosch, à coller depuis « Données à copier ».}\par
\vspace{3pt}
\begin{tcolorbox}[act]\textbf{M12 · Trois cales, trois questions} \hfill {\small 100 min · rotation sur postes · niveau 3 · mesures}\end{tcolorbox}
\textbf{Préparation.} Activité dense, pour un groupe à l'aise ou en fin de parcours. Le plan de rotation s'adapte au nombre de stagiaires dans la page (3 à 8).\par
\textbf{Animation.} Corrigé détaillé : \fichier{Guide\_Trois\_cales\_formateur.pdf}.\par
\textbf{\color{ok}Corrigé de la fiche M12.} Voir le guide dédié : chaque réponse y renvoie à la question de la fiche.\par
{\small\textbf{Repli ou remarque.} Bouton « Exemple simulé » de la page.}\par
\vspace{3pt}
\begin{tcolorbox}[act]\textbf{M13 · Bon ou mauvais : sommes-nous d'accord ?} \hfill {\small 45 min · en groupe · niveau 3 · Ellistat · mesures}\end{tcolorbox}
\textbf{Préparation.} Établir la référence au micromètre avant la séance. Sans butées proches des limites, tout le monde sera d'accord et l'activité n'apprend rien.\par
\textbf{Animation.} Le pied à coulisse est volontairement trop grossier pour cette tolérance (voir M1) : les désaccords viendront des pièces limites. Distinguer deux erreurs : accepter une mauvaise pièce (le client paie) et rejeter une bonne (l'atelier paie).\par
\textbf{\color{ok}Corrigé de la fiche M13.} Résultats propres au groupe. \textbf{Q3} Les désaccords se concentrent sur les butées proches de 11,85 et de 12,00 : la résolution de 0,02 mm ne permet pas de trancher.\par Le guide Ellistat donne un seul seuil : kappa au moins égal à 0,75. Kappa = 1 : accord parfait ; kappa = 0 : accord dû au seul hasard.\par
{\small\textbf{Repli ou remarque.} Activité à préparer : elle dépend des butées disponibles.}\par
\vspace{3pt}
\section*{Bloc A · Cartes aux attributs}
\begin{tcolorbox}[act]\textbf{A1 · Défectueuse ou défauts ?} \hfill {\small 15 min · en binôme · niveau 1}\end{tcolorbox}
\textbf{Animation.} Partir d'exemples des entreprises du groupe : que comptez-vous au poste ? La décision à prendre choisit le comptage : trier des pièces (p) ou améliorer un procédé (c).\par
\textbf{\color{ok}Corrigé de la fiche A1.} 1 : défectueuses, fixe, np (ou p). 2 : défauts, fixe, c. 3 : défectueuses, variable, p. 4 : défauts, variable, u. 5 : défauts, fixe, c. 6 : défectueuses, fixe, np (ou p).\par \textbf{Q1} 1 sur la carte p ; 3 sur la carte c. \textbf{Q2} np : pièces défectueuses et lots de taille fixe. c : défauts et unité de contrôle constante.\par
{\small\textbf{Repli ou remarque.} Ellistat propose p, Np, c, u et deux cartes de Laney (P', U').}\par
\vspace{3pt}
\begin{tcolorbox}[act]\textbf{A2 · Ma première carte p, à la main} \hfill {\small 35 min · en binôme · niveau 1}\end{tcolorbox}
\textbf{Préparation.} Donner le \textbf{scénario A} (stable) et un numéro de poste différent à chaque binôme : chacun aura ses propres lots.\par
\textbf{Animation.} Un lot à la fois : le rythme du poste de contrôle fait partie de l'activité. La carte de la page reste masquée jusqu'à la question Q4.\par
\textbf{\color{ok}Corrigé de la fiche A2.} Résultats propres à chaque poste : la page affiche $\bar p$, les limites et les points hors limites (volet « Afficher la carte »).\par Repère théorique (taux caché de 4 \%, $n = 50$) : limite supérieure $4 + 3 \times 2{,}77 = 12{,}3$ \%, soit 6,2 pièces ; limite inférieure 0. \textbf{Q3} Non : 5 pièces reste sous la limite.\par Jeu de secours (lots 1 à 20) : 38 défectueuses sur 1000, $\bar p = 3{,}8$ \% ; limite supérieure 11,9 \%, soit 5,96 pièces ; aucun point hors limites.\par Variante aux dés : $\bar p$ proche de 16,7 \%, limite supérieure vers 52 \% (5,2 dés).\par
{\small\textbf{Repli ou remarque.} Sans connexion : jeu « 35 lots de 50 pièces », lots 1 à 20.}\par
\vspace{3pt}
\begin{tcolorbox}[act]\textbf{A3 · La même carte dans Ellistat} \hfill {\small 20 min · en binôme · niveau 1 · Ellistat}\end{tcolorbox}
\textbf{Préparation.} Dérouler une fois l'écran des cartes aux attributs sur votre compte : le guide Ellistat n'en donne pas le pas à pas.\par
\textbf{Animation.} Si le collage échoue, changer le séparateur ou saisir les 20 valeurs à la main : 2 minutes.\par
\textbf{\color{ok}Corrigé de la fiche A3.} Les limites doivent être les mêmes que celles de A2, aux arrondis près. \textbf{Q2} La carte Np porte le nombre de pièces au lieu du pourcentage : même forme, mêmes alarmes. Au poste, compter des pièces est plus direct.\par
{\small\textbf{Repli ou remarque.} Jeu « 35 lots de 50 pièces », lots 1 à 20 : $\bar p = 3{,}8$ \%, limite supérieure 11,9 \%.}\par
\vspace{3pt}
\begin{tcolorbox}[act]\textbf{A4 · Le procédé se dérègle} \hfill {\small 30 min · en binôme · niveau 2 · Ellistat}\end{tcolorbox}
\textbf{Préparation.} Donner le \textbf{scénario B} (saut de 4 à 10 \% au lot 21) sans le dire. Pour un groupe à l'aise : scénario C (dérive lente) ou E (un seul lot anormal, le 18).\par
\textbf{Animation.} Ne rien annoncer : les stagiaires surveillent, la carte parle. Bouton « Révéler le taux caché » en bas de la page, à la fin seulement.\par
\textbf{\color{ok}Corrigé de la fiche A4.} Résultats propres à chaque poste. Avec un saut de 4 à 10 \% et $n = 50$, un lot dépasse la limite environ une fois sur cinq : l'alarme arrive en moyenne au bout de 4 à 5 lots. La suite de points au-dessus de la moyenne se voit avant.\par Jeu de secours : limites sur les lots 1 à 20 ($\bar p = 3{,}8$ \%, limite 11,9 \%) : lots 21, 25, 26 et 29 hors limites. Limites sur les 35 lots ($\bar p = 5{,}9$ \%, limite 15,9 \%) : plus aucun point ne sort.\par Variante aux dés : le taux passe de 16,7 à 33,3 \%. Avec 10 dés, la limite est à 5,2 : il faut 6 dés « défectueux ». L'alarme est lente (un lancer sur 13 environ) : petits lots, carte peu sensible.\par
{\small\textbf{Repli ou remarque.} Jeu « 35 lots de 50 pièces » : lots 1 à 20 à 4 \%, lots 21 à 35 à 10 \%.}\par
\vspace{3pt}
\begin{tcolorbox}[act]\textbf{A5 · p ou np ? Et si les lots changent de taille ?} \hfill {\small 25 min · en binôme · niveau 2 · Ellistat}\end{tcolorbox}
\textbf{Animation.} La prédiction se règle par une division : 4,5 \% contre 2,25 \%.\par
\textbf{\color{ok}Corrigé de la fiche A5.} $\bar p = 191 / 5970 = 3{,}20$ \%. Lot de 80 : limite $3{,}20 + 3 \times 1{,}97 = 9{,}10$ \%. Lot de 400 : $3{,}20 + 3 \times 0{,}88 = 5{,}84$ \%.\par \textbf{Q2} L'écart-type de la proportion vaut $\sqrt{\bar p(1-\bar p)/n}$ : il diminue quand $n$ augmente. \textbf{Q3} Le lot 18 : 27 défectueuses sur 300, soit 9,0 \% pour une limite de 6,25 \%. \textbf{Q4} 9 pièces sur 200 et 9 sur 400 ne disent pas la même chose : un nombre brut ne se compare qu'à taille égale.\par
{\small\textbf{Repli ou remarque.} Données simulées (taux de 3 \%, lot 18 dégradé volontairement).}\par
\vspace{3pt}
\begin{tcolorbox}[act]\textbf{A6 · Compter les défauts : la carte c} \hfill {\small 30 min · en binôme · niveau 2 · Ellistat}\end{tcolorbox}
\textbf{Préparation.} Scénario A pour les 20 premières unités. Pour la suite : scénario B (de 3 à 6 défauts au n° 21) ou C (dérive).\par
\textbf{Animation.} Faire nommer des défauts réels : rayure, bulle, grain, manque de matière.\par
\textbf{\color{ok}Corrigé de la fiche A6.} Résultats propres à chaque poste. Repère théorique (3 défauts en moyenne) : limite supérieure $3 + 3 \times 1{,}73 = 8{,}2$ ; limite inférieure 0. \textbf{Q3} 8 défauts reste sous la limite : rare, mais pas anormal.\par Jeu de secours (30 panneaux, limites sur les 20 premiers) : $\bar c = 83 / 20 = 4{,}15$ ; limite supérieure 10,26 ; seul le panneau 30 (11 défauts) sort. Les panneaux 21 à 30 sont tirés avec 8 défauts en moyenne : la carte c réagit lentement.\par \textbf{Q4} La taille de l'unité contrôlée : même surface, même longueur, même nombre de pièces.\par
{\small\textbf{Repli ou remarque.} Jeu « Défauts d'aspect sur 30 panneaux ».}\par
\vspace{3pt}
\begin{tcolorbox}[act]\textbf{A7 · Des unités de taille variable : la carte u} \hfill {\small 25 min · en binôme · niveau 3 · Ellistat}\end{tcolorbox}
\textbf{Animation.} Prédiction : 0,125 défaut par mètre contre 0,058. Le rouleau court est deux fois plus mauvais.\par
\textbf{\color{ok}Corrigé de la fiche A7.} $\bar u = 100 / 1820 = 0{,}055$ défaut par mètre. Rouleau 14 : 13 défauts sur 80 m, $u = 0{,}1625$ ; limite $0{,}055 + 3\sqrt{0{,}055 / 80} = 0{,}134$ : hors limites.\par \textbf{Q2} Seul le rouleau 14 sort. \textbf{Q3} La carte c compte les défauts sur des unités identiques ; la carte u les rapporte à la taille.\par
{\small\textbf{Repli ou remarque.} Données simulées (0,05 défaut par mètre, rouleau 14 dégradé volontairement).}\par
\vspace{3pt}
\begin{tcolorbox}[act]\textbf{A8 · Quelle période de référence ?} \hfill {\small 25 min · en binôme · niveau 2 · Ellistat}\end{tcolorbox}
\textbf{Animation.} Faire débattre : qui décide de la période de référence au poste ?\par
\textbf{\color{ok}Corrigé de la fiche A8.} Les 120 lots : $\bar p = 851 / 18\,000 = 4{,}73$ \% ; limite 9,93 \% ; lots 59, 68 et 112 hors limites (15, 16 et 17 non-conformes). Lots 1 à 20 : $\bar p = 3{,}43$ \% ; limite 7,89 \% ; 17 lots hors limites.\par \textbf{Q2} Non : plus calme que la suite. \textbf{Q3} Trop d'alarmes : on règle sans raison et on perd confiance dans la carte. Pas assez : on laisse passer des dérives.\par \textbf{Laney} $\sigma_z = 1{,}37$ ; limite P' 11,86 \% ; aucun lot ne sort (calcul de référence en Python, à comparer à l'écran).\par
{\small\textbf{Repli ou remarque.} Calculs refaits en Python sur le fichier.}\par
\vspace{3pt}
\begin{tcolorbox}[act]\textbf{A9 · La taille du lot compte} \hfill {\small 20 min · en binôme · niveau 2}\end{tcolorbox}
\textbf{Animation.} Activité courte, sans Ellistat : bonne transition entre A4 et A8.\par
\textbf{\color{ok}Corrigé de la fiche A9.} Repères théoriques (taux caché de 4 \%) : $n = 50$ : limites 0 et 12,3 \% ; $n = 500$ : limites 1,4 \% et 6,6 \%. Les valeurs relevées varient autour.\par \textbf{Q1} Environ par 3 ($\sqrt{10} = 3{,}2$). \textbf{Q2} Une amélioration réelle : à comprendre et à conserver. \textbf{Q3} Dès le lot 21 presque à coup sûr : 10 \% est très au-dessus de 6,6 \%.\par
{\small\textbf{Repli ou remarque.} Bouton « Révéler le taux caché » pour vérifier le scénario.}\par
\vspace{3pt}
\section*{Bloc C · Cartes aux mesures et capabilité}
\begin{tcolorbox}[act]\textbf{C1 · La carte sonne-t-elle ?} \hfill {\small 30 min · en binôme · niveau 1}\end{tcolorbox}
\textbf{Animation.} Imposer l'explication orale avant le clic : sans elle, la page se fait trop vite. Circuler et demander « quelle ligne regardes-tu ? » : la confusion limites et tolérances se voit là.\par
\textbf{\color{ok}Corrigé de la fiche C1.} La page corrige chaque carte. Règles et réactions : 1 point hors limites : plan de réaction du poste ; 7 points du même côté : chercher pourquoi le réglage a bougé, puis recentrer ; 6 points qui montent ou descendent : prévoir l'intervention (usure, échauffement) ; carte R au-dessus de sa limite : chercher ce qui a changé la dispersion ; aucune règle : ne toucher à rien.\par \textbf{Q2} Non : elles sont calculées à partir des mesures. \textbf{Q3} Régler sans raison ajoute de la dispersion.\par
{\small\textbf{Repli ou remarque.} Les trois pièges : point proche d'une limite, 5 points qui montent seulement, 6 points du même côté seulement. Aucun ne déclenche de règle.}\par
\vspace{3pt}
\begin{tcolorbox}[act]\textbf{C2 · La machine est-elle capable ?} \hfill {\small 30 min · en binôme · niveau 1}\end{tcolorbox}
\textbf{Animation.} Commencer par la série « donnés » : la décision d'abord, le calcul ensuite.\par
\textbf{\color{ok}Corrigé de la fiche C2.} La page corrige chaque cas. \textbf{Q1} La dispersion est bonne mais le procédé est décentré : on recentre. \textbf{Q2} La voiture est plus large que le garage : même centrée, des pièces sortent. Il faut réduire la dispersion (machine, outil, méthode). \textbf{Q3} Non : c'est Cpk qui dit ce qui arrivera sur les prochaines pièces.\par
{\small\textbf{Repli ou remarque.} Seuil de la page : Cpk au moins égal à 1,33.}\par
\vspace{3pt}
\begin{tcolorbox}[act]\textbf{C3 · La voiture et le garage} \hfill {\small 30 min · en binôme · niveau 2}\end{tcolorbox}
\textbf{Animation.} Laisser manipuler avant toute formule. L'onglet 2 est plus exigeant : pour un groupe à l'aise.\par
\textbf{\color{ok}Corrigé de la fiche C3.} \textbf{Défi 1} Par exemple σ = 0,018 (Cp = 1,39) et $m$ = 11,90 (Cpk = 0,93). \textbf{Défi 2} $m$ = 11,925 : Cpk = Cp = 1,16. Régler ne suffit pas. \textbf{Défi 3} σ = 0,15 / (6 × 1,33) = 0,0188. \textbf{Défi 4} $m$ = TI + 2,33 σ, par exemple σ = 0,02 et $m$ = 11,897.\par \textbf{Q1} Cpk seul. \textbf{Q3} Une moyenne de 5 pièces varie 2,24 fois moins qu'une pièce ($\sqrt 5$).\par
{\small\textbf{Repli ou remarque.} La page vérifie elle-même chaque défi.}\par
\vspace{3pt}
\begin{tcolorbox}[act]\textbf{C4 · Ma carte, à la main puis dans Ellistat} \hfill {\small 60 min · en binôme · niveau 2 · Ellistat}\end{tcolorbox}
\textbf{Animation.} Faire annoncer les limites au tableau avant d'ouvrir Ellistat.\par
\textbf{\color{ok}Corrigé de la fiche C4.} Moyennes : 11,864 ; 11,862 ; 11,884 ; 11,898 ; 11,910 ; 11,902. Étendues : 0,04 ; 0,05 ; 0,05 ; 0,04 ; 0,04 ; 0,08. $\bar{\bar X} = 11{,}8867$ ; $\bar R = 0{,}050$.\par À la main : limites 11,9155 et 11,8578 ; $LCS_R = 0{,}1057$ ; écart-type 0,0215 ; Cp = 1,16 ; Cpk = 0,57. Ellistat (écart-type groupé 0,0205) : limites 11,9142 et 11,8592 ; Cp = 1,22 ; Cpk = 0,60 ; Pp = 0,95 ; Ppk = 0,46.\par \textbf{Q1} Aucun point hors limites dans les deux cas. \textbf{Q2} L'écart-type groupé d'Ellistat. \textbf{Q3} La moyenne (11,887) est sous le milieu de la tolérance (11,925) : décentrage vers la limite basse. \textbf{Q4} Les moyennes montent de 11,862 à 11,910 : une dérive, que Pp voit et que Cp ignore.\par
{\small\textbf{Repli ou remarque.} Conventions d'Ellistat établies le 10/11/2026 sur une copie d'écran (50 butées) et recalculées en Python.}\par
\vspace{3pt}
\begin{tcolorbox}[act]\textbf{C5 · Le mauvais découpage} \hfill {\small 20 min · en binôme · niveau 3 · Ellistat}\end{tcolorbox}
\textbf{Animation.} Pour un groupe à l'aise. Bien faire dire ce que contient chaque ligne avant de calculer.\par
\textbf{\color{ok}Corrigé de la fiche C5.} Bon découpage : Cp = 1,22 ; Pp = 0,95. Mauvais découpage : Cp = 0,89 (écart-type groupé 0,0281) ; Pp = 0,95, inchangé.\par \textbf{Q2} Chaque ligne mélange des pièces de six moments différents : la dérive entre prélèvements gonfle la dispersion « court terme ». \textbf{Q3} Celui qui prélève.\par
{\small\textbf{Repli ou remarque.} Valeurs recalculées en Python avec l'écart-type groupé.}\par
\vspace{3pt}
\begin{tcolorbox}[act]\textbf{C6 · Cinquante butées qui dérivent} \hfill {\small 30 min · en binôme · niveau 2 · Ellistat}\end{tcolorbox}
\textbf{Animation.} Faire répondre à la prédiction par un vote avant d'ouvrir la carte.\par
\textbf{\color{ok}Corrigé de la fiche C6.} Cp = 1,34 ; Cpk = 1,05 ; Pp = 0,70 ; Ppk = 0,54 ; Rs = 0,52 ; Rr = 0,78 (valeurs lues dans Ellistat le 10/11/2026).\par \textbf{Q1} Sous-groupes 1 et 2 sous la limite basse (11,8836), 8, 9 et 10 au-dessus de la limite haute (11,9336). Les moyennes montent régulièrement. \textbf{Q2} Usure d'outil, dérive thermique, réglage qui glisse.\par \textbf{Q3} Un procédé hors contrôle n'est pas prévisible : rien ne dit que demain ressemblera à aujourd'hui. \textbf{Q4} Rs = Pp / Cp : la moitié du potentiel est perdue par instabilité. Rr = Ppk / Pp : perte due au décentrage.\par
{\small\textbf{Repli ou remarque.} Ellistat estime 5,1 \% de la population sous la limite basse.}\par
\vspace{3pt}
\begin{tcolorbox}[act]\textbf{C7 · Trois enquêtes} \hfill {\small 60 min · en binôme · niveau 2 · Ellistat}\end{tcolorbox}
\textbf{Préparation.} Répartir les enquêtes : Pression et Dimension pour tous les groupes, Liquide en plus pour les plus rapides.\par
\textbf{Animation.} Présentation de 5 minutes par binôme. Le formateur relie chaque cas au plan de réaction.\par
\textbf{\color{ok}Corrigé de la fiche C7.} \textbf{Pression} (limites Ellistat sur 24 jours : 63,16 et 51,44) : jours 15, 16 et 19 hors limites. À la main avec $A_2$ : 62,73 et 51,87, hors limites 15, 16, 18, 19, 23. Sur les 12 premiers jours ($A_2$) : 62,88 et 52,78, hors limites 15 à 20 et 23.\par \textbf{Dimension} (limites Ellistat : 2670,7 et 2169,0) : hors limites 2, 12, 20, 21, 22, 24, 25, 26 ; étendue du sous-groupe 2 hors limite. La valeur 1421 est aberrante : à vérifier d'abord. Du sous-groupe 19 au 27, tout est au-dessus de la moyenne : changement durable.\par \textbf{Liquide} : carte des valeurs individuelles sur les 15 premiers lots : limites 44,60 et 32,06 ; limite des étendues mobiles 7,70. Le lot 17 (44) reste dans les limites, mais son saut de 10 dépasse 7,70.\par \textbf{Q3} Non : les périodes perturbées élargissent les limites. On les calcule sur une période stable, puis on les fige.\par
{\small\textbf{Repli ou remarque.} Les alertes affichées dépendent des règles d'Ellistat (nombre de points non documenté) : relever ce que l'écran signale.}\par
\vspace{3pt}
\begin{tcolorbox}[act]\textbf{C8 · Du moyen de mesure à la décision} \hfill {\small 90 min · en binôme · niveau 3 · Ellistat}\end{tcolorbox}
\textbf{Préparation.} Les données de justesse et de R\&R sont simulées.\par
\textbf{Animation.} Possible en évaluation : noter l'ordre suivi, la justesse des lectures et la décision.\par
\textbf{\color{ok}Corrigé de la fiche C8.} \textbf{Justesse} : moyenne 12,00076 ; écart-type 0,000716 ; biais +0,00076. Formules du livret Bosch : Cg environ 7,0 et Cgk environ 6,6. Instrument capable.\par \textbf{R\&R} : EV 0,00108 ; AV 0,00090 ; GRR 0,00140 ; \%GRR 5,6 \% avec le facteur 6 (4,8 \% avec 5,15) ; ndc 34. Capable dans les deux conventions.\par \textbf{Production} : voir C6. \textbf{Décision} : la mesure est validée, donc les points hors limites viennent du procédé. Le procédé n'est pas sous contrôle : aucune capabilité à annoncer. Action : supprimer la dérive, recentrer, puis refaire l'étude.\par
{\small\textbf{Repli ou remarque.} Valeurs Ellistat de Cg et du \%GRR à confirmer (conventions non confirmées).}\par
\vspace{3pt}

\vspace{6pt}
{\footnotesize\color{uimm}\textbf{Sources.} Ellistat, guide utilisateur de Data Analysis (rubriques Général, Justesse, R\&R sur mesures, Cartes de contrôle aux attributs ; pages consultées le 11/11/2026) ; réponse d'Ellistat du 10/11/2026 sur les sous-groupes. Robert Bosch GmbH, \emph{Booklet 10 : Capability of Measurement and Test Processes}, édition novembre 2019 (§ 4.1, 4.2, 4.4, fig. 3, exemples p. 11 et p. 14). Tous les corrigés chiffrés ont été recalculés en Python (NumPy, SciPy). Les jeux marqués « simulé » sont tirés avec une graine fixe (1125, 2027, 2611).}
\end{document}
